\chapter{Problem Formulation}\label{chap:problem}

\section{The Primal Problem (P)}

We begin with the soft margin classification problem, also known as the \emph{primal problem}, which is a convex optimization problem:
For ease of notation, we denote the problem in a more compact vectorized form:

\begin{align*}\label{eq:primal-problem}
  \min_{\symbf{w}, b, \symbf{\xi}} \quad & f(\symbf{w}, b, \symbf{\xi}) = \frac{1}{2}\|\symbf{w}\|_2^2 + C\symbf{1}_M^\top \symbf{\xi} \text{subject to} \quad &
  \begin{cases}
    \symbf{1}_M - Y(X\symbf{w} + b\symbf{1}_M) - \symbf{\xi} & \leq \symbf{0}_M \\
    \symbf{\xi}                                              & \geq \symbf{0}_M
  \end{cases}\tag{P}
\end{align*}

Here, $\symbf{w} \in \R^d$ represents the weight vector, $b \in \R$ is the bias term, and $\symbf{\xi} \in \R^M$ is the vector of slack variables, which account for misclassified points by allowing constraint violations.

The training data consists of $M$ data points, represented by the matrix $X = [\symbf{x}_1, \ldots, \symbf{x}_M] \in \R^{d \times M}$,
where each column $\symbf{x}_i$ is a data point associated with a class label $y_i \in {\pm 1}$, collectively forming the label vector $\symbf{y}=[y_1, \ldots, y_M]^T \in \R^M$.

For convenience, we define the diagonal class-label matrix $Y = \operatorname{diag}(y_1, \ldots, y_M) \in \R^{M \times M}$, where $y_i$ the class label for the $i$-th data point.

The parameter $C > 0$ controls the trade-off between maximizing the margin and minimizing classification error. A higher $C$ penalizes misclassification more heavily, enforcing stricter adherence to the margin constraints.

\begin{theorem}{Existence of Solution for Primal Problem}{primal-existence}
  The primal problem~\eqref{eq:primal-problem} has atleast one solution (i.e. a global minimizer) $(\symbf{w}^\star, b^\star, \symbf{\xi}^\star)$ if the feasible set is non-empty, and the objective function $f$ is coercive.
\end{theorem}

\begin{proof}{}{}
  The primal problem~\eqref{eq:primal-problem} is a convex quadratic program with linear constraints.
  To establish existence of a solution, we need to verify that the feasible set is non-empty, and the objective function is coercive.

  \paragraph{Feasibility}
  Let $\symbf{w} = \symbf{0}_d$, $b = 0$, and $\symbf{\xi} = \symbf{1}_M$. 
  This point satisfies all constraints since:
  \begin{align*}
    \symbf{1}_M - Y(X\symbf{w} + b\symbf{1}_M) - \symbf{\xi} = \symbf{0}_M \leq \symbf{0}_M \quad \text{and} \quad \symbf{\xi} = \symbf{1}_M \geq \symbf{0}_M
  \end{align*}

  \paragraph{Coercivity}  
  The dominating term $\frac{1}{2}\norm{\symbf{w}}_2^2 \to \infty$ when $\symbf{w} \to \infty$, ensuring that the objective function is coercive.

  
  Since the feasible set is non-empty and closed, and $f$ is coercive, the \emph{Weierstrass theorem} guarantees existence of a global minimizer\cite{weierstrass}.\qed

\end{proof}

\begin{theorem}{Uniqueness of Solution for Primal Problem}{primal-uniqueness}
  The weight vector $\symbf{w}^\star$ is unique.

  The offset $b^\star$ is unique if there exists at least one support vector $\symbf{x}_i$ with $0 < \alpha_i^\star < C$.

  The slack variables $\symbf{\xi}^\star$ are not necessarily unique.
\end{theorem}
\begin{proof}{}{}
  \paragraph{Uniqueness of $\symbf{w}^\star$}
  The objective is strictly convex in $\symbf{w}$ due to the $\|\symbf{w}\|_2^2$ term. 
  Strict convexity ensures a unique minimizer for $\symbf{w}$.

  \paragraph{Uniqueness of $b^\star$}
  From the KKT conditions, for any $i$ with $0 < \alpha_i^\star < C$ (margin support vectors):
  \[
    y_i(\inner{\symbf{w}^\star, \symbf{x}_i} + b^\star) = 1 \implies b^\star = y_i - \inner{\symbf{w}^\star, \symbf{x}_i}
  \]
  If such $i$ exist, $b^\star$ is uniquely determined. If not, $b^\star$ may lie in an interval.

  \paragraph{Non-Uniqueness of $\symbf{\xi}^\star$}
  For $i$ with $\alpha_i^\star = C$, $\xi_i = \max\{0, 1-y_i(\inner{\symbf{w}^\star, \symbf{x}_i} + b^\star)\}$ is fixed.
  For $0 < \alpha_i^\star < C$, $\xi_i = 0$.
  For $\alpha_i^\star = 0$, $\xi_i$ is unconstrained but bounded by $\xi_i \geq 0$.
\end{proof}

\section{The Dual Problem (DP)}

The dual of the primal problem~\eqref{eq:primal-problem} can be derived using Lagrange multipliers, leading to the following dual optimization problem:

\begin{align}\label{eq:dual-problem}
  \min_{\symbf{\alpha}}\frac{1}{2} \inner*{\symbf{\alpha}, \mathrm{Y} \mathrm{G} \mathrm{Y} \symbf{\alpha}} - \inner*{\symbf{1}_M, \symbf{\alpha}}
  \quad \text{subject to} \quad
  h(\symbf{\alpha}, \symbf{y}) =
  \begin{cases}
    \inner*{\symbf{y}, \symbf{\alpha}} = 0 \\
    \symbf{0}_M \leq \symbf{\alpha} \leq C\symbf{1}_M
  \end{cases}
\end{align}

where \(\mathrm{G} = \left(\inner*{\symbf{x}_i, \symbf{x}_j}\right)_{i,j=1}^M\) is the Gram matrix of the points.

\begin{theorem}{Existence and Uniqueness of Solution for Dual Problem}{}
 
\end{theorem}

\begin{proof}{}{}
\end{proof}

